\chapter{Space groups and their representations}\label{ch:space}

\begin{watch}{\lec{6}}
Complex irreps and bilinear invariants \ts{6}{3}{0:03} $\cdot$
$\Cv{4}$ versus $C_4$ \ts{6}{314}{5:14} $\cdot$
circular polarisation and time reversal \ts{6}{1081}{18:01} $\cdot$
from atoms to crystals; the bands of graphene \ts{6}{1519}{25:19} $\cdot$
symmetry-enforced vs accidental degeneracies \ts{6}{1775}{29:35} $\cdot$
irreps of the translation group \ts{6}{2724}{45:24} $\cdot$
the Brillouin zone as the set of irreps \ts{6}{3232}{53:52} $\cdot$
Bloch's theorem as Wigner's theorem \ts{6}{3295}{54:55} $\cdot$
Bloch sums from orbitals \ts{6}{3832}{63:52} $\cdot$
space groups, symmorphic and non-symmorphic \ts{6}{4217}{70:17} $\cdot$
star and little group \ts{6}{4661}{77:41} $\cdot$
small representations \ts{6}{5175}{86:15}
\end{watch}

\section{Complex representations and bilinear invariants}\label{sec:complexinv}

So far all irreps were real. When an irrep is complex, the invariant scalar product
of two copies must be the \emph{hermitian} one, $\sum_iz_i^{*}w_i$, because the matrices
are unitary rather than orthogonal. Two consequences:
\begin{itemize}
\item a bilinear invariant pairs an irrep with its complex conjugate, as in
      \eqref{eq:ninv};
\item $z^{*}w$ is complex: its real and imaginary parts are two independent real
      invariants.
\end{itemize}

Compare $\Cv{4}$ with its rotation subgroup $C_4=\{E,C_4^+,C_2,C_4^-\}$ (rotations about
$z$ only). $C_4$ is Abelian, every element is a class, so there are four irreps, all
one-dimensional (the only solution of $\sum d_\mu^2=4$ with four terms):
\begin{equation}\label{eq:C4table}
\begin{array}{c|cccc|l}
C_4 & E & C_4^+ & C_2 & C_4^- & \text{components}\\\hline
\irr{A} & 1&1&1&1 & E_z\\
\irr{B} & 1&-1&1&-1 & \\
{}^{1}\irr{E} & 1&\ii&-1&-\ii & E_+=E_x+\ii E_y\\
{}^{2}\irr{E} & 1&-\ii&-1&\ii & E_-=E_x-\ii E_y
\end{array}
\end{equation}
\begin{quote}\small
\emph{Convention.} Here we list vector \emph{components}: the rotation $C_4^+$ maps
$(E_x,E_y)\mapsto(-E_y,E_x)$, hence $E_+\mapsto\ii E_+$. If instead one lists
\emph{basis functions} with $(\hat{O}_gf)(\vec{r})=f(g^{-1}\vec{r})$, the function $x+\ii y$
picks up $-\ii$. The two conventions are complex conjugates of each other; tables
differ on this point, so check before using one.
\end{quote}
Every irrep of an Abelian group is one-dimensional. This will be crucial for the
translation group.

Physically, $E_\pm$ are circular polarisations. A rotation by $90^\circ$ shifts their
phase by $\pm90^\circ$. Time reversal reverses the sense of rotation and exchanges
them. The pair ${}^{1}\irr{E}\oplus{}^{2}\irr{E}$ is therefore glued by time reversal
into a two-dimensional \emph{physically irreducible} representation, which is why it
carries the letter $\irr{E}$.

\paragraph{Invariants.} In $\Cv{4}$, $\GV=\irr{A}_1\oplus\irr{E}$ as in $\Cv{3}$, and
the bilinear invariants of two vectors $\vec{E}$, $\vec{P}$ are $E_zP_z$ and
$E_xP_x+E_yP_y$. In $C_4$, $\GV=\irr{A}\oplus{}^{1}\irr{E}\oplus{}^{2}\irr{E}$, and
\[
E_+^{*}P_+=(E_xP_x+E_yP_y)+\ii\,(E_xP_y-E_yP_x)
\]
provides, besides $E_zP_z$, \emph{two} invariants: the in-plane scalar product and
$(\vec{E}\times\vec{P})_z$. The second is allowed because $C_4$ has no mirrors.
Groups made only of proper rotations are called \emph{chiral}; they allow such
``vector-product'' couplings, with physical consequences such as optical activity.

\paragraph{Checking symmetry-adapted coordinates.} To verify that given coordinates
belong to an irrep, transform them with the \emph{generators} only (for $C_4$, just
$C_4^+$; for $\Cv{4}$, $C_4^+$ and one mirror); every other element is a product of
generators. For a multidimensional irrep compare with the tabulated matrices, since
the character is not enough.

\section{From molecules to crystals}

A crystal is a ``molecule'' (a set of atoms, the basis) repeated at the points of a
Bravais lattice
\[
\vec{t}=n_1\vec{a}_1+n_2\vec{a}_2+n_3\vec{a}_3,\qquad n_i\in\mathbb{Z}.
\]
For graphene the basis consists of two carbon atoms, A and B, identical but with
different environments (\cref{fig:honeycomb}). The crystal is invariant under all
lattice translations, which form the Abelian \emph{translation group} $\Tgroup$.
Taken literally, a perfect crystal has infinitely many degrees of freedom and its
representations are infinite dimensional. These are the \emph{band representations}
of \cref{ch:topology}. The translation group lets us break them into finite pieces.

\begin{figure}[ht]
\centering
\begin{tikzpicture}[scale=0.95]
  \def\s{1.7320508}
  % hexagon centres at t = n1 a1 + n2 a2 with a1=(1,0), a2=(1/2,sqrt3/2) (scale a=1.6)
  \begin{scope}[scale=1.6]
  \foreach \n in {-2,...,2}{\foreach \m in {-2,...,2}{
    \pgfmathsetmacro{\cx}{\n+0.5*\m}\pgfmathsetmacro{\cy}{0.8660254*\m}
    \pgfmathtruncatemacro{\inside}{(abs(\cx)<2.3) && (abs(\cy)<1.6) ? 1 : 0}
    \ifnum\inside=1
      \foreach \k in {0,...,5}{
        \draw[gray!60] ($(\cx,\cy)+(30+60*\k:0.57735)$)--($(\cx,\cy)+(90+60*\k:0.57735)$);}
      \fill[red!70!black] ($(\cx,\cy)+(30:0.57735)$) circle (0.045);
      \fill[red!70!black] ($(\cx,\cy)+(150:0.57735)$) circle (0.045);
      \fill[red!70!black] ($(\cx,\cy)+(270:0.57735)$) circle (0.045);
      \fill[green!50!black] ($(\cx,\cy)+(90:0.57735)$) circle (0.045);
      \fill[green!50!black] ($(\cx,\cy)+(210:0.57735)$) circle (0.045);
      \fill[green!50!black] ($(\cx,\cy)+(330:0.57735)$) circle (0.045);
    \fi}}
  \draw[-{Stealth},thick] (0,0)--(1,0) node[below right]{$\vec{a}_1$};
  \draw[-{Stealth},thick] (0,0)--(0.5,0.8660254) node[above]{$\vec{a}_2$};
  \fill (0,0) circle (0.035) node[below left,font=\scriptsize]{$O$};
  \node[red!70!black,font=\scriptsize,right] at ($(30:0.57735)+(0.03,0)$) {A};
  \node[green!50!black,font=\scriptsize,above] at ($(90:0.57735)+(0,0.02)$) {B};
  \end{scope}
\end{tikzpicture}
\caption{The honeycomb lattice. The origin $O$ is the centre of a hexagon (Wyckoff
position $1a$, site symmetry $\Cv{6}$). The atoms of sublattice A (red) sit at
$\vec{q}_{\mathrm{A}}=\tfrac13(\vec{a}_1+\vec{a}_2)$ and those of sublattice B (green) at
$\vec{q}_{\mathrm{B}}=\tfrac23(\vec{a}_1+\vec{a}_2)$ plus lattice vectors (Wyckoff position
$2b$, site symmetry $\Cv{3}$). Lattice constant $a=\abs{\vec{a}_1}$.}
\label{fig:honeycomb}
\end{figure}

\section{Irreducible representations of the translation group}

Since $\Tgroup$ is Abelian, its irreps are one-dimensional: to each translation a
unimodular complex number (a phase), with $D(\vec{t}_1)D(\vec{t}_2)=D(\vec{t}_1+\vec{t}_2)$.
The only continuous functions turning sums into products are exponentials, so
\begin{equation}\label{eq:Tirrep}
D^{\vec{k}}(\vec{t})=\ee^{-\ii\vec{k}\cdot\vec{t}}
\end{equation}
for some wave vector $\vec{k}$. (Imposing Born--von K\'arm\'an periodic boundary
conditions on a crystal of $N$ cells makes $\vec{k}$ discrete, with $N$ allowed values;
nothing else changes.)

Introduce the reciprocal basis, $\vec{a}_i\cdot\vec{b}_j=2\cpi\,\delta_{ij}$, and write
$\vec{k}=\sum_ik_i\vec{b}_i$. Then $\vec{k}\cdot\vec{t}=2\cpi\sum_ik_in_i$. If
$\vec{G}=\sum_im_i\vec{b}_i$ with integer $m_i$ (a reciprocal-lattice vector),
$\ee^{-\ii\vec{G}\cdot\vec{t}}=1$ for every $\vec{t}$. Hence $\vec{k}$ and $\vec{k}+\vec{G}$
label the same irrep: they are \emph{equivalent}, written $\vec{k}\equiv\vec{k}+\vec{G}$.

\begin{keyresult}[The Brillouin zone]
The inequivalent irreps of $\Tgroup$ are in one-to-one correspondence with the points
of any primitive cell of the reciprocal lattice. The first Brillouin zone is the
most symmetric such cell. \emph{The Brillouin zone is the set of irreps of the
translation group}, and $\vec{k}$ (``crystal momentum'') is a label of an irrep, just as
$l$ labels an irrep of $\SO{3}$.
\end{keyresult}

\section{Bloch's theorem is Wigner's theorem}

Let the translation operators act on functions as $(\hat{T}_{\vec{t}}\psi)(\vec{r})=
\psi(\vec{r}-\vec{t})$ (active convention: the function is moved by $+\vec{t}$). A
function belongs to the irrep $\vec{k}$ if
$\hat{T}_{\vec{t}}\psi=\ee^{-\ii\vec{k}\cdot\vec{t}}\psi$, i.e.\
$\psi(\vec{r}+\vec{t})=\ee^{\ii\vec{k}\cdot\vec{t}}\psi(\vec{r})$. The general solution is
\begin{equation}
\psi_{\vec{k}}(\vec{r})=\ee^{\ii\vec{k}\cdot\vec{r}}\,u_{\vec{k}}(\vec{r}),\qquad
u_{\vec{k}}(\vec{r}+\vec{t})=u_{\vec{k}}(\vec{r}).
\end{equation}
Since the Hamiltonian of a crystal commutes with $\Tgroup$, its eigenfunctions can be
chosen to belong to irreps of $\Tgroup$ (Wigner). That is all Bloch's theorem says.
(Bradley and Cracknell use the passive convention, with the opposite sign of the
exponent in \eqref{eq:Tirrep}.)

\subsection{Bloch sums}
Place an orbital $\varphi_\alpha$ at position $\vec{q}_\alpha$ of every cell and form
\begin{equation}\label{eq:blochsum}
\Phi_{\alpha\vec{k}}(\vec{r})=\sum_{\vec{t}}\ee^{\ii\vec{k}\cdot\vec{t}}\,
\varphi_\alpha(\vec{r}-\vec{t}-\vec{q}_\alpha).
\end{equation}
Relabelling the sum shows $\hat{T}_{\vec{t}'}\Phi_{\alpha\vec{k}}=\ee^{-\ii\vec{k}\cdot\vec{t}'}
\Phi_{\alpha\vec{k}}$: the Bloch sum belongs to the irrep $\vec{k}$. With this
convention $\Phi_{\alpha,\vec{k}+\vec{G}}=\Phi_{\alpha\vec{k}}$. Bloch sums are
``Bloch-like'' waves: eigenfunctions of translations but not in general of $\hat{H}$.

Because $\hat{H}$ is invariant, it has no matrix elements between inequivalent irreps
of $\Tgroup$ (\cref{sec:countinv}). With $N_{\mathrm{o}}$ orbitals per cell, the
$\infty\times\infty$ tight-binding problem in real space splits into one
$N_{\mathrm{o}}\times N_{\mathrm{o}}$ problem at each $\vec{k}$, giving $N_{\mathrm{o}}$ bands.
Different points of the Brillouin zone do not talk to each other. This is the only
use of symmetry you can make without group theory. To go further we need the rest of
the space group.

\section{Space groups}

A space-group element combines a point operation $R$ (rotation, reflection,
rotoinversion) with a translation $\vec{v}$. In Seitz notation
\begin{equation}
\seitz{R}{\vec{v}}\,\vec{r}=R\vec{r}+\vec{v},
\end{equation}
first $R$, then the translation. From this,
\begin{equation}\label{eq:seitz}
\seitz{R_1}{\vec{v}_1}\seitz{R_2}{\vec{v}_2}=\seitz{R_1R_2}{\vec{v}_1+R_1\vec{v}_2},\qquad
\seitz{R}{\vec{v}}^{-1}=\seitz{R^{-1}}{-R^{-1}\vec{v}}.
\end{equation}
The pure translations $\seitz{E}{\vec{t}}$ form a normal subgroup $\Tgroup$, and the set
of rotational parts forms the \emph{point group} $\bar{G}\cong G/\Tgroup$
(\cref{sec:normal,sec:semidirect}).

\begin{definition}
A space group is \emph{symmorphic} if, for a suitable origin, every element has the
form $\seitz{R}{\vec{t}}$ with $\vec{t}\in\Tgroup$; equivalently, $G$ contains the point
group $\bar{G}$ as a subgroup, $G=\Tgroup\rtimes\bar{G}$. Otherwise it is
\emph{non-symmorphic}: some operations come only with a fractional translation that
no choice of origin removes. These are \emph{glide planes} (reflection plus a half
translation parallel to the plane) and \emph{screw axes} (rotation $C_n$ plus a
translation $\tfrac{p}{n}\vec{t}$ along the axis).
\end{definition}
In three dimensions there are 230 space groups, of which 73 are symmorphic. In two
dimensions there are 17 plane (wallpaper) groups, 13 of them symmorphic; the
honeycomb lattice has the symmorphic plane group $p6mm$. Its three-dimensional
analogue, used in the lecture and in the Bilbao tables, is $P6mm$ (No.~183). For
symmorphic groups the representation theory reduces to that of point groups; for
non-symmorphic groups it is subtler but fully tabulated.

\subsection{Semidirect products\beyondmark}\label{sec:semidirect}
The translation subgroup $\Tgroup$ is normal (\cref{sec:normal}):
$\seitz{R}{\vec{v}}\seitz{E}{\vec{t}}\seitz{R}{\vec{v}}^{-1}=\seitz{E}{R\vec{t}}$, and
$R\vec{t}$ is again a lattice vector (\cref{ex:6:seitz}). The quotient $G/\Tgroup$ is
isomorphic to the point group $\bar{G}$: the coset of $\seitz{R}{\vec{v}}$ is determined
by $R$ alone.

If $G$ has a normal subgroup $N$ and a subgroup $H$ with $N\cap H=\{e\}$ and $NH=G$,
then $G=N\rtimes H$ is a \emph{semidirect product}: $H$ acts on $N$ by conjugation.
(When $H$ is normal too, this is the direct product of \cref{sec:directproduct}.)
Examples: $\Cv{3}=C_3\rtimes\{E,\sd\}$; a symmorphic space group is
$\Tgroup\rtimes\bar{G}$, the point group acting on the lattice by rotation.

Non-symmorphic space groups are \emph{not} semidirect products: $\bar{G}$ is still the
quotient $G/\Tgroup$ but is not a subgroup of $G$, because some rotations only occur
combined with fractional translations. This is the structural reason why their
representation theory needs projective representations (see the box on non-symmorphic
groups in \cref{sec:star}).

\subsection{Space-group operations move \texorpdfstring{$\vec{k}$}{k}}
Let $g=\seitz{R}{\vec{v}}$ act on functions by $(\hat{O}_g\psi)(\vec{r})=\psi(g^{-1}\vec{r})$.
From \eqref{eq:seitz}, $g^{-1}\seitz{E}{\vec{t}}g=\seitz{E}{R^{-1}\vec{t}}$. Hence, if
$\psi$ belongs to $\vec{k}$,
\[
\hat{T}_{\vec{t}}\,\hat{O}_g\psi=\hat{O}_g\hat{T}_{R^{-1}\vec{t}}\psi
=\ee^{-\ii\vec{k}\cdot R^{-1}\vec{t}}\,\hat{O}_g\psi
=\ee^{-\ii(R\vec{k})\cdot\vec{t}}\,\hat{O}_g\psi :
\]
\emph{$\hat{O}_g$ maps the irrep $\vec{k}$ of $\Tgroup$ onto the irrep $R\vec{k}$.}

\section{Star, little group and small representations}\label{sec:star}

\begin{definition}
\begin{itemize}
\item The \emph{star} of $\vec{k}$ is the set of inequivalent vectors
      $\{R\vec{k}\,:\,R\in\bar{G}\}$.
\item The \emph{little group} $G_{\vec{k}}$ is the set of $g=\seitz{R}{\vec{v}}\in G$
      with $R\vec{k}\equiv\vec{k}$, i.e.\ $R\vec{k}=\vec{k}+\vec{G}$ for some reciprocal
      vector $\vec{G}$. Its point parts form the \emph{little co-group}
      $\bar{G}_{\vec{k}}\subset\bar{G}$.
\item A \emph{small representation} is an irrep of $G_{\vec{k}}$ in which the
      translations are represented by $\ee^{-\ii\vec{k}\cdot\vec{t}}\,\one$.
\end{itemize}
\end{definition}
The equivalence ``up to $\vec{G}$'' is essential: for graphene the corner $K$ of the
hexagonal zone is carried by $C_3$ to another corner that differs from it by a
reciprocal-lattice vector, so $C_3$ belongs to the little co-group of $K$
(\cref{ch:bands}).

Why the little group matters: the Bloch states at a given $\vec{k}$ are mapped onto
themselves by $G_{\vec{k}}$, and the $N_{\mathrm{o}}\times N_{\mathrm{o}}$ Hamiltonian
$\mat{H}(\vec{k})$ commutes with the small representation. By Wigner's theorem the
energy levels at $\vec{k}$ are labelled by small irreps, and their degeneracies are the
dimensions of those irreps. For bilinear quantities (Hamiltonians, optical matrix
elements between states at the same $\vec{k}$) only small representations are needed.
Full space-group irreps are needed for electron--phonon processes, superstructures
and similar problems. They are obtained by \emph{induction} from small irreps and
have dimension (size of the star)$\times$(dimension of the small irrep).

\begin{keyresult}[Small irreps of symmorphic groups]
If $G$ is symmorphic, the small irreps at $\vec{k}$ are
\begin{equation}\label{eq:smallsym}
D^{\vec{k},\mu}\bigl(\seitz{R}{\vec{t}}\bigr)=\ee^{-\ii\vec{k}\cdot\vec{t}}\,
\mat{D}^{(\mu)}(R),\qquad R\in\bar{G}_{\vec{k}},\ \vec{t}\in\Tgroup,
\end{equation}
where $\mat{D}^{(\mu)}$ runs over the irreps of the little co-group (a point group).
They are labelled like point-group irreps, or as $\Gamma_1,\Gamma_2,\dots$,
$K_1,K_2,\dots$ in the Bilbao convention.
\end{keyresult}

\begin{beyond}[Beyond the lectures: non-symmorphic groups]
For $g=\seitz{R}{\vec{v}}\in G_{\vec{k}}$ with fractional $\vec{v}$, try
$\mat{\Delta}(g)=\ee^{-\ii\vec{k}\cdot\vec{v}}\mat{D}(R)$. Since
$g_1g_2=\seitz{R_1R_2}{\vec{v}_1+R_1\vec{v}_2}$, the product
$\mat{\Delta}(g_1)\mat{\Delta}(g_2)=\ee^{-\ii\vec{k}\cdot(\vec{v}_1+\vec{v}_2)}\mat{D}(R_1)\mat{D}(R_2)$
must equal $\ee^{-\ii\vec{k}\cdot(\vec{v}_1+R_1\vec{v}_2)}\mat{D}(R_1R_2)$. The two phases
differ by $\ee^{-\ii\vec{k}\cdot\vec{v}_2+\ii\vec{k}\cdot R_1\vec{v}_2}=
\ee^{\ii(R_1^{-1}\vec{k}-\vec{k})\cdot\vec{v}_2}$. If every $R\in\bar{G}_{\vec{k}}$ fixes
$\vec{k}$ exactly (always the case inside the zone), this factor is $1$ and $\mat{D}$
is an ordinary irrep of $\bar{G}_{\vec{k}}$. On the zone boundary,
$R_1^{-1}\vec{k}=\vec{k}+\vec{G}_1$ and the factor $\ee^{\ii\vec{G}_1\cdot\vec{v}_2}$ can
differ from $1$ when $\vec{v}_2$ is fractional. Then $\mat{D}$ must be a
\emph{projective} representation of the little co-group. This is how glides and screws
force bands to stick together on the zone boundary (\cref{ex:6:glide}).
\end{beyond}

\section{Degeneracies: enforced or accidental?}\label{sec:accidental}

Look at a first-principles band structure of graphene (\cref{fig:grapheneTB} shows the
$\uppi$ bands). Its bands carry labels of small irreps at $\Gamma$, $K$, $M$ and along the
lines joining them. Two kinds of degeneracy appear:
\begin{itemize}
\item \emph{Symmetry-enforced}: at $K$ two bands meet because they carry a
      two-dimensional small irrep. Any Hamiltonian with the symmetry of graphene must
      have a degeneracy there: the Dirac point.
\item \emph{Accidental}: along a symmetry line two bands with \emph{different}
      one-dimensional irreps cross. Group theory permits the crossing (the
      Hamiltonian cannot couple different irreps) but does not require it.
      Moving the bands changes where they cross, or removes the crossing.
\end{itemize}
Two further observations. At an isolated high-symmetry point an accidental
degeneracy requires fine-tuning of the energies to infinite precision, so it has
probability zero. If you find one numerically, either the precision is insufficient
or you have missed a symmetry. Along a line, $\vec{k}$ is a continuous parameter, and
crossings of different irreps happen generically. Two bands of the \emph{same} irrep
generically repel (avoided crossing), because they can be coupled.

Near a degeneracy the bands disperse linearly unless a symmetry forbids the linear
term. Group theory also decides this, through the method of invariants
(\cref{sec:kp}).

\section{Time reversal in crystals\beyondmark}\label{sec:TRcrystal}

\Cref{sec:TR} treated time reversal for an atom or a molecule. In a crystal it also
acts on the wave vector.

\subsection{Time reversal and Bloch states}
Complex conjugation turns $\ee^{\ii\vec{k}\cdot\vec{r}}u_{\vec{k}}$ into
$\ee^{-\ii\vec{k}\cdot\vec{r}}u^{*}_{\vec{k}}$: time reversal maps $\vec{k}$ to $-\vec{k}$.
Hence
\begin{itemize}
\item $E_n(\vec{k})=E_n(-\vec{k})$ (without spin), or $E_{n\uparrow}(\vec{k})=E_{n\downarrow}(-\vec{k})$
      and in general Kramers pairs at $\vec{k}$ and $-\vec{k}$ (with spin);
\item at the \emph{time-reversal-invariant momenta} (TRIM), $\vec{k}\equiv-\vec{k}$ (in
      graphene: $\Gamma$ and the three $M$ points), time reversal maps the states at
      $\vec{k}$ onto themselves. With $\TR^2=-1$ every band is then Kramers degenerate;
      in the presence of inversion symmetry as well, every band is doubly degenerate
      at every $\vec{k}$.
\end{itemize}
$K$ and $K'=-K$ in graphene are exchanged by time reversal; $K$ is not a TRIM.

At a TRIM of a symmorphic space group, the test of \cref{sec:herring} applies to the
little co-group. For general $\vec{k}$ and non-symmorphic groups the same question is
answered by Herring's criterion, a sum of $\chi\bigl((\TR g)^2\bigr)$ over the elements
$g$ that map $\vec{k}$ to $-\vec{k}$; see Bradley--Cracknell.

\subsection{Magnetic groups}\label{sec:magnetic}
In a magnetically ordered crystal time reversal alone is broken, but a combination
$\TR g$ with a spatial operation $g$ may survive. For example, an antiferromagnet can be
invariant under $\TR$ combined with a translation or with inversion. The groups that
contain such antiunitary elements are the \emph{magnetic} (Shubnikov) groups: 1651
magnetic space groups in three dimensions, of four types (no $\TR$ at all; $\TR$ itself,
the ``grey'' groups; $\TR$ only combined with point operations; $\TR$ combined with a
translation). Their representation theory uses Wigner's \emph{corepresentations}, whose
degeneracies are again decided by tests of the type above. The Bilbao server tabulates
their irreducible corepresentations and band representations (magnetic topological
quantum chemistry).

\exercisesection

\begin{exercise}\label{ex:6:C4v}
Show that $(\vec{E}\times\vec{P})_z=E_xP_y-E_yP_x$ is invariant under $C_4$ but changes
sign under the mirrors of $\Cv{4}$. To which irrep of $\Cv{4}$ does it belong?
\end{exercise}

\begin{solution}
Under $C_4^+$, $(E_x,E_y)\to(-E_y,E_x)$ and likewise for $\vec{P}$:
$E_x'P_y'-E_y'P_x'=(-E_y)P_x-E_x(-P_y)=E_xP_y-E_yP_x$. Under the mirror
$x\to-x$: $(-E_x)P_y-E_y(-P_x)=-(E_xP_y-E_yP_x)$. So it is invariant under the rotations
and odd under the mirrors: $\irr{A}_2$ of $\Cv{4}$, like $R_z$.
\end{solution}

\begin{exercise}\label{ex:6:recip}
For the hexagonal lattice $\vec{a}_1=a(1,0)$, $\vec{a}_2=a(\tfrac12,\tfrac{\sqrt3}{2})$:
(a) find $\vec{b}_1,\vec{b}_2$; (b) show that $K=\tfrac23\vec{b}_1+\tfrac13\vec{b}_2$ is a
corner of the Brillouin zone with $\abs{K}=4\cpi/(3a)$, and that $M=\tfrac12\vec{b}_1$ is
the midpoint of an edge; (c) show that the six corners fall into two classes of three
equivalent points, $K$ and $K'\equiv-K$, and that $K\not\equiv K'$.
\end{exercise}

\begin{solution}
(a) $\vec{b}_1=\tfrac{2\cpi}{a}(1,-\tfrac{1}{\sqrt3})$, $\vec{b}_2=\tfrac{2\cpi}{a}(0,\tfrac{2}{\sqrt3})$,
$\abs{\vec{b}_i}=\tfrac{4\cpi}{\sqrt3a}$.
(b) $K=\tfrac23\vec{b}_1+\tfrac13\vec{b}_2=(\tfrac{4\cpi}{3a},0)$. The zone is a hexagon with
inradius $\abs{\vec{b}}/2=\tfrac{2\cpi}{\sqrt3a}$ and circumradius
$\tfrac{2\cpi}{\sqrt3a}/\cos30^\circ=\tfrac{4\cpi}{3a}=\abs{K}$: $K$ is a corner. $M=\vec{b}_1/2$
has $\abs{M}$ equal to the inradius and lies on the edge bisector: the midpoint of an
edge.
(c) $C_3^+K-K=\tfrac{4\cpi}{3a}(-\tfrac32,\tfrac{\sqrt3}{2})=-\vec{b}_1$, so $C_3K\equiv K$; the
corners at $0^\circ,120^\circ,240^\circ$ are equivalent, and so are those at
$60^\circ,180^\circ,300^\circ$ (equivalent to $K'=-K$). But $K'-K=-2K=(-\tfrac{8\cpi}{3a},0)$
would need $m_1=-\tfrac43$ in $m_1\vec{b}_1+m_2\vec{b}_2$: not a reciprocal vector, so
$K\not\equiv K'$.
\end{solution}

\begin{exercise}\label{ex:6:seitz}
Prove the composition and inverse rules \eqref{eq:seitz}. Show that
$g\seitz{E}{\vec{t}}g^{-1}=\seitz{E}{R\vec{t}}$ for $g=\seitz{R}{\vec{v}}$, so that
$\Tgroup$ is a normal subgroup.
\end{exercise}

\begin{solution}
$\seitz{R_1}{\vec{v}_1}\seitz{R_2}{\vec{v}_2}\vec{r}=R_1(R_2\vec{r}+\vec{v}_2)+\vec{v}_1
=R_1R_2\vec{r}+\vec{v}_1+R_1\vec{v}_2$.
$\seitz{R^{-1}}{-R^{-1}\vec{v}}\seitz{R}{\vec{v}}\vec{r}=R^{-1}(R\vec{r}+\vec{v})-R^{-1}\vec{v}=\vec{r}$.
Finally
$\seitz{R}{\vec{v}}\seitz{E}{\vec{t}}\seitz{R^{-1}}{-R^{-1}\vec{v}}
=\seitz{R}{\vec{v}+R\vec{t}}\seitz{R^{-1}}{-R^{-1}\vec{v}}=\seitz{E}{R\vec{t}}$, and $R\vec{t}$ is a
lattice vector because $R$ maps the lattice onto itself.
\end{solution}

\begin{exercise}\label{ex:6:bloch}
Verify from \eqref{eq:blochsum} that $\hat{T}_{\vec{t}'}\Phi_{\alpha\vec{k}}=
\ee^{-\ii\vec{k}\cdot\vec{t}'}\Phi_{\alpha\vec{k}}$ and $\Phi_{\alpha,\vec{k}+\vec{G}}=
\Phi_{\alpha\vec{k}}$. What changes if the phase $\ee^{\ii\vec{k}\cdot(\vec{t}+\vec{q}_\alpha)}$
is used instead?
\end{exercise}

\begin{solution}
$\hat{T}_{\vec{t}'}\Phi_{\alpha\vec{k}}(\vec{r})=\sum_{\vec{t}}\ee^{\ii\vec{k}\cdot\vec{t}}
\varphi_\alpha(\vec{r}-\vec{t}'-\vec{t}-\vec{q}_\alpha)=\sum_{\vec{s}}\ee^{\ii\vec{k}\cdot(\vec{s}-\vec{t}')}
\varphi_\alpha(\vec{r}-\vec{s}-\vec{q}_\alpha)=\ee^{-\ii\vec{k}\cdot\vec{t}'}\Phi_{\alpha\vec{k}}$
with $\vec{s}=\vec{t}+\vec{t}'$. Since $\ee^{\ii\vec{G}\cdot\vec{t}}=1$, $\Phi_{\alpha,\vec{k}+\vec{G}}=\Phi_{\alpha\vec{k}}$.
With the phase $\ee^{\ii\vec{k}\cdot(\vec{t}+\vec{q}_\alpha)}$ each Bloch sum is multiplied by
$\ee^{\ii\vec{k}\cdot\vec{q}_\alpha}$: the translation property is unchanged, but now
$\Phi_{\alpha,\vec{k}+\vec{G}}=\ee^{\ii\vec{G}\cdot\vec{q}_\alpha}\Phi_{\alpha\vec{k}}$, so $\mat{H}(\vec{k})$ is
periodic only up to a diagonal unitary transformation. Both conventions are in use.
The one with $\vec{q}_\alpha$ makes $\mat{H}(\vec{k})$ reflect the positions of the orbitals
(useful for Berry phases); the other makes it periodic.
\end{solution}

\begin{exercise}[a glide line]\label{ex:6:glide}
The plane group $pg$ (rectangular lattice $a\vec{e}_x$, $b\vec{e}_y$) is generated by
translations and the glide $g=\seitz{m_x}{\tfrac{b}{2}\vec{e}_y}$, where
$m_x\colon(x,y)\mapsto(-x,y)$. (a) Show that $g^2=\seitz{E}{b\vec{e}_y}$. (b) On the line
$k_x=0$, Bloch states can be chosen as eigenstates of $\hat{O}_g$; find the possible
eigenvalues as functions of $k_y$. (c) Follow an eigenvalue continuously from
$k_y=-\cpi/b$ to $k_y=\cpi/b$ and show that the two eigenvalue branches are exchanged:
bands on this line come in connected pairs.
\end{exercise}

\begin{solution}
(a) $g^2=\seitz{m_x^2}{\tfrac{b}{2}\vec{e}_y+m_x\tfrac{b}{2}\vec{e}_y}=\seitz{E}{b\vec{e}_y}$, since
$m_x$ leaves $\vec{e}_y$ unchanged.
(b) On $k_x=0$, $m_x\vec{k}=\vec{k}$, so $g$ is in the little group and $\hat{O}_g$ can be
diagonalised together with $\hat{H}$. Its eigenvalues satisfy
$\lambda^2=\ee^{-\ii k_yb}$ (the eigenvalue of $\hat{T}_{b\vec{e}_y}$), so
$\lambda_\pm(k_y)=\pm\ee^{-\ii k_yb/2}$.
(c) $\lambda_+(-\cpi/b)=\ee^{\ii\cpi/2}=\ii$ and $\lambda_+(\cpi/b)=\ee^{-\ii\cpi/2}=-\ii=\lambda_-(-\cpi/b)$.
Since $k_y=\cpi/b$ and $k_y=-\cpi/b$ are the same point, a band that starts in the
branch $\lambda_+$ at $-\cpi/b$ arrives in the branch $\lambda_-$: it cannot close on itself
and must connect to a partner. Bands along this line come in connected pairs, a
consequence of the fractional translation.
\end{solution}
